% GENERATED FILE. Edit canonical lessons, then run npm run generate:books.
% canonical-source-hash: fnv1a64:acf1af6459ae37d9
% canonical-lessons: LA-C331-luctus, LA-C331-nenia, LA-C331-laudatio, LA-C331-monumentum, LA-C331-inferiae

\chapter{Lūctus}
\label{ch:la-luctus}
\hlchaptermodality{331}

\begin{chapteropening}
\textbf{By the end of this chapter:} \emph{I can talk about mourning and remembrance: grief, the dirge, the speech of praise, the memorial and the offerings to the dead.}

Complete the last lesson of chapter 331: I can talk about mourning and remembrance: grief, the dirge, the speech of praise, the memorial and the offerings to the dead.
\end{chapteropening}

\section[lūctus, lūctūs]{lūctus, lūctūs — grief, mourning}

\label{lesson:LA-C331-luctus}



\begin{quote}
Before the new one: say the Latin for a funeral pyre, then the Latin for a funeral procession.
\end{quote}

\subsection*{You'll want to know: lūctus, lūctūs}
\textbf{lūctus, lūctūs} — \textquotedblleft{}grief, mourning\textquotedblright{}.

\begin{quote}
Tōta domus in lūctū erat.
\end{quote}

\begin{quote}
\emph{The whole household was in mourning.}
\end{quote}

Masculine, 4th declension, from \emph{lūgeō}, 'to mourn' (taught). Mourners wore dark clothes.

\begin{cousinweb}
\textbf{lūctus} is built on \emph{lūgeō}, \textquotedblleft{}to mourn\textquotedblright{} (EtymDB): sorrow shown for the dead.
\end{cousinweb}

\begin{grammarlens}[title={What you've built}]
One more word for telling, explaining and describing at B1.
\end{grammarlens}

\subsection*{Your turn}
\begin{itemize}
\raggedright
  \item \emph{Say it:} \emph{lūctus, lūctūs}
  \item \emph{Say it:} \emph{lūctus, lūctūs}, once more
  \item \emph{Say it:} a sentence of your own with \emph{lūctus, lūctūs}
  \item \emph{Recall:} say the Latin for a funeral pyre, then the Latin for a funeral procession, then say \emph{lūctus, lūctūs} again
\end{itemize}

\begin{tcolorbox}[breakable,colback=teal!4,colframe=teal!35!black,title={Before you move on}]
What does lūctus, lūctūs mean? (\textquotedblleft{}Grief, mourning\textquotedblright{}.) Say it once more.
\end{tcolorbox}
\section[nēnia, nēniae]{nēnia, nēniae — a dirge, a funeral song}

\label{lesson:LA-C331-nenia}



\begin{quote}
Before the new one: say the Latin for a funeral procession, then the Latin for grief.
\end{quote}

\subsection*{You'll want to know: nēnia, nēniae}
\textbf{nēnia, nēniae} — \textquotedblleft{}a dirge, a funeral song\textquotedblright{}.

\begin{quote}
Fēminae nēniam trīstem canēbant.
\end{quote}

\begin{quote}
\emph{The women were singing a sad dirge.}
\end{quote}

Feminine, also spelled \emph{naenia}. Hired women sang it in praise of the dead; later also a lullaby or a silly song.

\begin{cousinweb}
\textbf{nēnia} is a word whose origin is not settled (Lewis \& Short).
\end{cousinweb}

\begin{grammarlens}[title={What you've built}]
One more word for telling, explaining and describing at B1.
\end{grammarlens}

\subsection*{Your turn}
\begin{itemize}
\raggedright
  \item \emph{Say it:} \emph{nēnia, nēniae}
  \item \emph{Say it:} \emph{nēnia, nēniae}, once more
  \item \emph{Say it:} a sentence of your own with \emph{nēnia, nēniae}
  \item \emph{Recall:} say the Latin for a funeral procession, then the Latin for grief, then say \emph{nēnia, nēniae} again
\end{itemize}

\begin{tcolorbox}[breakable,colback=teal!4,colframe=teal!35!black,title={Before you move on}]
What does nēnia, nēniae mean? (\textquotedblleft{}A dirge, a funeral song\textquotedblright{}.) Say it once more.
\end{tcolorbox}
\section[laudātiō, laudātiōnis]{laudātiō, laudātiōnis — a speech of praise, a eulogy}

\label{lesson:LA-C331-laudatio}



\begin{quote}
Before the new one: say the Latin for grief, then the Latin for a dirge.
\end{quote}

\subsection*{You'll want to know: laudātiō, laudātiōnis}
\textbf{laudātiō, laudātiōnis} — \textquotedblleft{}a speech of praise, a eulogy\textquotedblright{}.

\begin{quote}
Fīlius in forō laudātiōnem patris habuit.
\end{quote}

\begin{quote}
\emph{The son gave a speech in praise of his father in the forum.}
\end{quote}

Feminine, from \emph{laudō}, 'to praise' (taught). \emph{Laudātiō fūnebris}: the funeral speech.

\begin{cousinweb}
\textbf{laudātiō} is built on \emph{laudō}, \textquotedblleft{}to praise\textquotedblright{} (Lewis \& Short), with the ending \emph{-tiō}.
\end{cousinweb}

\begin{grammarlens}[title={What you've built}]
One more word for telling, explaining and describing at B1.
\end{grammarlens}

\subsection*{Your turn}
\begin{itemize}
\raggedright
  \item \emph{Say it:} \emph{laudātiō, laudātiōnis}
  \item \emph{Say it:} \emph{laudātiō, laudātiōnis}, once more
  \item \emph{Say it:} a sentence of your own with \emph{laudātiō, laudātiōnis}
  \item \emph{Recall:} say the Latin for grief, then the Latin for a dirge, then say \emph{laudātiō, laudātiōnis} again
\end{itemize}

\begin{tcolorbox}[breakable,colback=teal!4,colframe=teal!35!black,title={Before you move on}]
What does laudātiō, laudātiōnis mean? (\textquotedblleft{}A speech of praise, a eulogy\textquotedblright{}.) Say it once more.
\end{tcolorbox}
\section[monumentum, monumentī]{monumentum, monumentī — a memorial, a monument}

\label{lesson:LA-C331-monumentum}



\begin{quote}
Before the new one: say the Latin for a dirge, then the Latin for a speech of praise.
\end{quote}

\subsection*{You'll want to know: monumentum, monumentī}
\textbf{monumentum, monumentī} — \textquotedblleft{}a memorial, a monument\textquotedblright{}.

\begin{quote}
Fīliī patrī monumentum ad viam posuērunt.
\end{quote}

\begin{quote}
\emph{The sons set up a memorial to their father by the road.}
\end{quote}

Neuter, from \emph{moneō}, 'to remind' (taught): anything that keeps a memory alive, a tomb, a statue or a book. English \emph{monument}.

\begin{cousinweb}
\textbf{monumentum} is built on \emph{moneō}, \textquotedblleft{}to remind\textquotedblright{}, with the ending \emph{-mentum} (EtymDB): a reminder. English \emph{monument} comes from it.
\end{cousinweb}

\begin{grammarlens}[title={What you've built}]
One more word for telling, explaining and describing at B1.
\end{grammarlens}

\subsection*{Your turn}
\begin{itemize}
\raggedright
  \item \emph{Say it:} \emph{monumentum, monumentī}
  \item \emph{Say it:} \emph{monumentum, monumentī}, once more
  \item \emph{Say it:} a sentence of your own with \emph{monumentum, monumentī}
  \item \emph{Recall:} say the Latin for a dirge, then the Latin for a speech of praise, then say \emph{monumentum, monumentī} again
\end{itemize}

\begin{tcolorbox}[breakable,colback=teal!4,colframe=teal!35!black,title={Before you move on}]
What does monumentum, monumentī mean? (\textquotedblleft{}A memorial, a monument\textquotedblright{}.) Say it once more.
\end{tcolorbox}
\section[īnferiae, īnferiārum]{īnferiae, īnferiārum — offerings to the dead}

\label{lesson:LA-C331-inferiae}



\begin{quote}
Before the new one: say the Latin for a speech of praise, then the Latin for a memorial.
\end{quote}

\subsection*{You'll want to know: īnferiae, īnferiārum}
\textbf{īnferiae, īnferiārum} — \textquotedblleft{}offerings to the dead\textquotedblright{}.

\begin{quote}
Fīliī īnferiās ad sepulcrum patris tulērunt.
\end{quote}

\begin{quote}
\emph{The sons brought offerings to their father's tomb.}
\end{quote}

Feminine, plural only. Wine, milk and flowers were poured or laid at the grave.

\begin{cousinweb}
\textbf{īnferiae} is built on \emph{īnferī}, \textquotedblleft{}those below, the dead\textquotedblright{}, from \emph{īnferus}, \textquotedblleft{}lower\textquotedblright{} (Lewis \& Short): gifts for those below.
\end{cousinweb}

\begin{grammarlens}[title={What you've built}]
One more word for telling, explaining and describing at B1.
\end{grammarlens}

\subsection*{Your turn}
\begin{itemize}
\raggedright
  \item \emph{Say it:} \emph{īnferiae, īnferiārum}
  \item \emph{Say it:} \emph{īnferiae, īnferiārum}, once more
  \item \emph{Say it:} a sentence of your own with \emph{īnferiae, īnferiārum}
  \item \emph{Recall:} say the Latin for a speech of praise, then the Latin for a memorial, then say \emph{īnferiae, īnferiārum} again
\end{itemize}

\begin{tcolorbox}[breakable,colback=teal!4,colframe=teal!35!black,title={Before you move on}]
What does īnferiae, īnferiārum mean? (\textquotedblleft{}Offerings to the dead\textquotedblright{}.) Say it once more.
\end{tcolorbox}
